\documentclass{amsart}
\usepackage{amssymb,amsmath,amsthm,hyperref}
\hypersetup{
  hidelinks,
  pdftitle={A density-one lower bound for the first decrease of a harmonic denominator},
  pdfauthor={Wouter van Doorn}
}

\setlength\parindent{0pt}
\hyphenpenalty=10000

\newtheorem{theorem}{Theorem}
\newtheorem{lemma}[theorem]{Lemma}
\newtheorem{corollary}[theorem]{Corollary}

\begin{document}

\title[A density-one lower bound]{A density-one lower bound for the first decrease of a harmonic denominator}

\author{ChatGPT}
\address{}
\email{}
\date{}

\begin{abstract}
For a positive integer $a$, let $b(a)$ be the smallest integer $b>a$ such that the denominator of
$$
\frac{1}{a}+\frac{1}{a+1}+\cdots+\frac{1}{b}
$$
is smaller than the denominator of the same sum ending at $b-1$. We prove that
$$
b(a)>a+\exp\left(\frac12\sqrt{\log a\log\log a}\right)
$$
for almost all positive integers $a$. In particular, for every fixed $C>0$ we have $b(a)>a+(\log a)^C$ for almost all $a$.
\end{abstract}

\maketitle

\section{Introduction}
For positive integers $a<b$, let $u_{a,b}$ and $v_{a,b}$ be the coprime positive integers for which
$$
\sum_{i=a}^b\frac1i=\frac{u_{a,b}}{v_{a,b}}.
$$
All logarithms in this paper are natural. For a given $a\in\mathbb N$, let $b(a)$ be the smallest integer $b>a$ for which $v_{a,b}<v_{a,b-1}$. The existence of $b(a)$ for every $a$ was asked by Erd\H{o}s and Graham in~\cite{cnt} and proved in~\cite{VD}. More recently, the exact positive value of
$$
\liminf_{a\to\infty}\frac{b(a)-a}{\log a}
$$
was determined in~\cite{shortest}. Thus there are infinitely many $a$ for which $b(a)-a$ has logarithmic order. The result below shows that these values of $a$ form a set of density zero.

\begin{theorem}\label{main}
For almost all positive integers $a$ we have
$$
b(a)>a+\exp\left(\frac12\sqrt{\log a\log\log a}\right).
$$
Equivalently, as $x\to\infty$,
$$
\#\left\{a\le x:b(a)\le a+\exp\left(\frac12\sqrt{\log a\log\log a}\right)\right\}=o(x).
$$
\end{theorem}

\begin{corollary}\label{powers}
For every fixed $C>0$, we have
$$
b(a)>a+(\log a)^C
$$
for almost all positive integers $a$. In particular, there are infinitely many positive integers $a$ for which $b(a)>a+(\log a)^2$.
\end{corollary}

\section{A necessary condition for a decrease}
For a positive integer $n$, let
$$
L_n:=\operatorname{lcm}(1,2,\ldots,n).
$$
We start with a necessary condition for the denominator to decrease.

\begin{lemma}\label{local}
For all positive integers $a<b$, we have
$$
\frac{v_{a,b}}{v_{a,b-1}}\ge \frac{b}{\gcd(b,L_{b-a})^2}.
$$
Consequently, if $v_{a,b}<v_{a,b-1}$, then
$$
\gcd(b,L_{b-a})>\sqrt b.
$$
\end{lemma}

\begin{proof}
Write
$$
\sum_{i=a}^{b-1}\frac1i=\frac UV,
$$
where $U$ and $V$ are coprime positive integers, so that $V=v_{a,b-1}$. Let $q=\gcd(V,b)$ and write $V=qV'$ and $b=qb'$. We then have $\gcd(V',b')=1$ and
$$
\frac UV+\frac1b=\frac{Ub'+V'}{qV'b'}.
$$
Since $\gcd(U,V')=\gcd(V',b')=1$, the numerator $Ub'+V'$ is coprime to both $V'$ and $b'$. Hence the only possible cancellation is with $q$, and
$$
\frac{v_{a,b}}{v_{a,b-1}}
=\frac{b'}{\gcd(Ub'+V',q)}
\ge \frac{b'}q
=\frac{b}{q^2}.
$$

Since $V$ divides $\operatorname{lcm}(a,a+1,\ldots,b-1)$, we have
$$
q\le \gcd\big(b,\operatorname{lcm}(a,a+1,\ldots,b-1)\big).
$$
For every prime power $p^k\mid b$, there is a multiple of $p^k$ in the interval $[a,b-1]$ if and only if $b-p^k\ge a$, or equivalently $p^k\le b-a$. It follows that
$$
\gcd\big(b,\operatorname{lcm}(a,a+1,\ldots,b-1)\big)=\gcd(b,L_{b-a}).
$$
This proves the first claim. If $v_{a,b}<v_{a,b-1}$, the first claim gives $b/\gcd(b,L_{b-a})^2<1$, which proves the second claim.
\end{proof}

\section{Proof of Theorem~\ref{main}}
\begin{proof}
Let $X$ be large and set
$$
N:=\left\lfloor\exp\left(\frac12\sqrt{\log(2X)\log\log(2X)}\right)\right\rfloor.
$$
We count the integers $a\in[X,2X)$ for which $v_{a,a+n}<v_{a,a+n-1}$ for some $1\le n\le N$. By Lemma~\ref{local}, every such pair $(a,n)$ satisfies
$$
\gcd(a+n,L_n)>\sqrt{a+n}>\sqrt X.
$$
Moreover, this greatest common divisor divides $L_N$ and $a+n$. Since $N<X$ for all sufficiently large $X$, it is also smaller than $3X$. Therefore the number of integers under consideration is at most
\begin{equation}\label{count}
N\sum_{\substack{d\mid L_N\\ \sqrt X<d<3X}}\left(\frac Xd+1\right).
\end{equation}

Let $0<\delta\le\frac12$ and define
$$
E:=\prod_{p\le N}\left(1-p^{-1+\delta}\right)^{-1}.
$$
Since every divisor of $L_N$ only has prime factors at most $N$, we have
$$
\sum_{d\mid L_N}d^{-1+\delta}\le E.
$$
Moreover, $d<3X$ implies $X/d+1<4X/d$. Hence, after dividing~\eqref{count} by $X$, the proportion of integers $a\in[X,2X)$ under consideration is
\begin{equation}\label{proportion}
\ll N\sum_{\substack{d\mid L_N\\d>\sqrt X}}\frac1d
\le NE X^{-\delta/2}.
\end{equation}

We next estimate $E$. Since $\delta\le\frac12$, we have
$$
\log E\ll\sum_{p\le N}p^{-1+\delta}.
$$
By partial summation and the prime number theorem,
$$
\sum_{p\le N}p^{-1+\delta}
\ll \frac{N^\delta}{\log N}+\int_2^N\frac{t^{-1+\delta}}{\log t}\,dt.
$$
If $\delta\log N\ge1$, splitting the integral at $e^{1/\delta}$ gives
\begin{equation}\label{eulerbound}
\log E\ll \log\left(\frac2\delta\right)+\frac{N^\delta}{\delta\log N}.
\end{equation}
Indeed, the part below $e^{1/\delta}$ is $\ll\log(2/\delta)$. After the substitution $u=\delta\log t$, the part above $e^{1/\delta}$ is bounded by a constant multiple of
$$
\int_1^{\delta\log N}\frac{e^u}{u}\,du
\ll \frac{N^\delta}{\delta\log N}.
$$

Put $y=\log X$ and choose
$$
\delta:=\frac{\frac12\log y+\frac14\log\log y}{\log N}.
$$
For sufficiently large $X$, this choice satisfies $0<\delta\le\frac12$ and $\delta\log N\ge1$. Moreover,
$$
N^\delta=y^{1/2}(\log y)^{1/4},
$$
so~\eqref{eulerbound} gives
\begin{equation}\label{Esimple}
\log E\ll \log y+y^{1/2}(\log y)^{-3/4}.
\end{equation}
On the other hand, since
$$
\log N=\frac12\sqrt{y\log y}+o(1),
$$
we have
\begin{align}
\frac{\delta y}{2}-\log N
&=\frac{y}{2\log N}\left(\frac12\log y+\frac14\log\log y\right)-\log N \notag\\
&=\left(\frac14+o(1)\right)\sqrt{\frac{y}{\log y}}\log\log y.\label{saving}
\end{align}
The right-hand side of~\eqref{Esimple} is
$$
o\left(\sqrt{\frac{y}{\log y}}\log\log y\right).
$$
Combining this with~\eqref{saving}, we obtain
$$
\log\left(NE X^{-\delta/2}\right)
=\log N+\log E-\frac{\delta y}{2}\longrightarrow-\infty.
$$
By~\eqref{proportion}, only $o(X)$ integers $a\in[X,2X)$ admit a decrease within the next $N$ terms.

The function
$$
\exp\left(\frac12\sqrt{\log x\log\log x}\right)
$$
is increasing for all sufficiently large $x$. Hence every $a\in[X,2X)$ for which
$$
b(a)\le a+\exp\left(\frac12\sqrt{\log a\log\log a}\right)
$$
is among the integers just counted. Thus the number of exceptions in $[X,2X)$ is $o(X)$. A standard dyadic decomposition now shows that the number of exceptions up to $x$ is $o(x)$, proving Theorem~\ref{main}.
\end{proof}

\begin{proof}[Proof of Corollary~\ref{powers}]
For every fixed $C>0$, we have
$$
\frac12\sqrt{\log a\log\log a}>C\log\log a
$$
for all sufficiently large $a$. The corollary follows from Theorem~\ref{main}.
\end{proof}

\begin{thebibliography}{9}

\bibitem{VD}
W. van Doorn,
\emph{On the non-monotonicity of the denominator of generalized harmonic sums},
arXiv:2411.03073 (2024).

\bibitem{shortest}
W. van Doorn,
\emph{The shortest harmonic sums with decreasing denominator},
arXiv:2609.00104 (2026).

\bibitem{cnt}
P.~Erd\H{o}s and R.~L. Graham,
\emph{Old and new problems and results in combinatorial number theory},
Monographies de L'Enseignement Math\'ematique, \textbf{28},
1980.

\end{thebibliography}

\end{document}
